\documentclass[10pt,a4paper]{amsart}
\usepackage[margin=19mm]{geometry}
\usepackage{amsmath,amssymb,mathtools}
\usepackage[scr=rsfs,cal=euler]{mathalfa}
\usepackage{xfrac}
\usepackage[unicode,hidelinks,bookmarks=false]{hyperref}
\newcommand{\ie}{i.e.\ }
\newcommand{\cf}{cf.\ }
\newcommand{\resp}{respectively\ }
\newcommand{\Subsection}{\subsection}
\providecommand{\nobreakdash}{\nobreak\hbox{-}}
\setlength{\emergencystretch}{2em}
\multlinegap0pt
\usepackage{xcolor}
\usepackage{caption}
\usepackage{amssymb}
\usepackage[shortlabels]{enumitem}
\newtheorem{prop}{Proposition}
\newcommand{\CP}{\mathbb{CP}}
\title{Topological Realization of Algebras of Quasi-invariants, I\\ (with an Appendix by M. V. Feigin and K. E. Feldman)}
\author{Yuri Berest, Ajay C. Ramadoss, M. V. Feigin and K. E. Feldman}
\date{Source: arXiv:2305.10604v2 (CC BY 4.0)}
\begin{document}
\maketitle

\noindent\emph{Source and licence.} Topological Realization of Algebras of Quasi-invariants, I\\ (with an Appendix by M. V. Feigin and K. E. Feldman). Version arXiv:2305.10604v2 (CC BY 4.0), distributed by arXiv under the Creative Commons Attribution 4.0 International licence, http://creativecommons.org/licenses/by/4.0/, as recorded in the arXiv licence metadata for this exact version and shown on the abstract page https://arxiv.org/abs/2305.10604v2. Full licence notice: \texttt{licenses/arxiv-2305.10604-CC-BY-4.0.txt}.\par\smallskip

\noindent\textit{Editorial scope.} Counted unit: the article's prop third (source0.tex lines 3465--3470). The article's complete original proof of it is delivered with it (source0.tex lines 3471--3496, one \texttt{proof} environment of 26 lines). No mathematics is written, repaired or re-derived: the statement and the proof are the article's own lines, in the article's own notation, and no retained line is substituted or re-flowed.  Retained uncounted as the source-local context the proof needs: lines 3341--3344.  Nothing was omitted from any retained span, so the package carries no omission record.  The retained lines cite 1 works, whose entries are transcribed verbatim from the article's own bibliography source in the e-print and delivered with short editorial labels recorded in the item's reference mapping.  No retained line contains a figure environment, an image-inclusion command or drawing code, so no image is delivered and the package declares no asset dependency.  Editorial formatting only, in addition: an AMS \texttt{amsart} preamble that restates the article's own declarations and macros used by the retained lines, namely the environments \texttt{prop} on separate counters, together with the standard packages those lines need.  The retained environments print in this document as Proposition 1 in order of appearance, whereas the article prints the same items with the numbering of its own full document.  The complete native source of the article as distributed in the arXiv e-print for this exact version is bundled unchanged as \texttt{sources/141674/source0.tex}.
\begin{equation}
    \label{formexpanded}
\prod_{j=1}^n(z_j x +w_j y)(\bar{w}_j x - \bar{z}_jy) = f_0x^{2n} + \cdots + f_{2n}y^{2n},
\end{equation}

\begin{prop}\label{third}
    The functions $\frac{f_0}{f_n},\frac{f_1}{f_n},\dots,\frac{f_{n-1}}{f_n}$ of the coefficients of the form \eqref{formexpanded} induce the structure of a complex vector space in the tangent spaces to $S^{2n}$ at the fixed points  $p_1$ and $p_2$ of the $T$-action. The representation $\chi\colon T \to {\rm End}(\mathbb C^n)$ of the torus $T$ in the tangent spaces of $S^{2n}$ at the points $p_1$ and $p_2$ with respect to this complex structure has weights 
    \[
    (-n,n), (-n+1, n-1), \ldots, (-1, 1). 
    \]
\end{prop}

    \begin{proof}
    
        Near the point $p_1\in S^{2n}$ we can assume that its preimage under the map $\widehat{Ar}$ 
        has a representative $\left((z_1,1), \ldots, (z_n, 1)\right)$, where 
%        satisfies $w_i=1$ for $1\le i \le n$, and 
$|z_i|$ is small. Let $\sigma_i$ be the $i$th elementary symmetric polynomial in $z_1, \ldots, z_n$, and let $\bar \sigma_i$ be its complex conjugate, which is the elementary symmetric polynomial in $\bar z_1, \ldots, \bar z_n$. We can take $\sigma_i, \bar\sigma_i$ as local coordinates near the preimage $\tilde q_1$ of $p_1$  under the map 
    $$
    G\colon \mathbb{C}P(1)^n/S_n\cong \mathbb{C}P(n) \to \mathbb{C}P(1)^n/W\cong S^{2n}.
    $$ Note that (cf. \cite{Arnold1996})
\begin{equation}
    \label{fsigma}
f_k = \sum_{j=0}^k (-1)^j \bar \sigma_j \sigma_{n-k+j}, \quad 0\le k \le n,
\end{equation}
and $f_{2n-k} = (-1)^{n-k} \bar f_k$. In particular, $f_n$ is real and $f_n(0)=1$. The latter implies that $f_i/f_n$ for  $0\le i \le n-1$ can be taken as complex local coordinates on $S^{2n}$ near the point $p_1$. 
 
The action of the torus on $S^{2n}$ is induced by the map $G$. Note that for $t=(t_1, t_2) \in T$ we have 
\begin{equation*}
\label{tz}   
t (z_i:1) = ( t_1 t_2^{-1} z_i:1).
\end{equation*}
  In the local coordinates $\sigma_i,\bar \sigma_i$ we have 
$t(\sigma _i) = (t_1 t_2^{-1})^i \sigma_i$ and $t(\bar \sigma_i) = (t_1 t_2^{-1})^{-i} \bar\sigma_i$. Therefore $t(f_k/f_n) = (t_1 t_2^{-1})^{n-k} f_k/f_n$ by formula \eqref{fsigma}, which implies the statement. 

The claim for $p_2$ follows similarly. In this case points near a point $\tilde q_2\in \CP(1)^n/S_n$ such that $G(\tilde q_2)=p_2$ can be represented as $\left((z_1, 1), \ldots, (z_{n-1}, 1), (1, z_n)\right)$, where $|z_i|$ is small for all $i$. 
Let $\sigma_i$ now be the $i$th elementary symmetric polynomial in $z_1, \ldots, z_{n-1}$ and $-\bar z_n$, and let the coordinate $\bar \sigma_i$ be its complex conjugate. Then  formula \eqref{fsigma} holds after multiplying the right-hand side by -1. The action of an element $t=(t_1, t_2)\in T$  on the coordinate $z_i$ is given by $t(z_i) = t_1 t_2^{-1} z_i$ for $1\le i \le n-1$, and $t(z_n)=t_1^{-1} t_2 z_n$, $t(\bar z_n)=t_1 t_2^{-1} \bar z_n$, which implies the statement. 
    \end{proof}

\begin{thebibliography}{99}
\bibitem[Arnold]{Arnold1996} V.~Arnold, \emph{Topological content of the {M}axwell theorem on multipole representation of spherical functions}, Topol. Methods Nonlinear Anal. \textbf{7} (1996), no.~2, 205--217.
\end{thebibliography}
\end{document}
